%%%%%%%%%%%%%%%%%%%%%%%%%%%%%%%%%%%%%%%%%%%%%%%%%%%%%%%%%%%%%%%%%%
% KAPITEL 5: DIE ALGEBRAISCHE STRUKTUR DER HALBGLATTEN SCHICHT
%%%%%%%%%%%%%%%%%%%%%%%%%%%%%%%%%%%%%%%%%%%%%%%%%%%%%%%%%%%%%%%%%%
\section{Die algebraische Struktur der halbglatten Schicht $\mathcal{S}_{\mathrm{halb}}$}
\label{sec:schalenstruktur}
Nachdem in den vorigen Abschnitten die rationale Struktur des Kollaps-Operators $\Pi_\Gamma$ über $\mathbb{Q}$ numerisch verifiziert wurde, gilt es nun, den arithmetischen Träger in der halbglatten Schicht $\mathcal{S}_{\mathrm{halb}} = \langle 5, 7 \rangle$ exakt zu bestimmen. Ein naiver Lift des Vierzustandsraums in den Kreisteilungskörper $\mathbb{Q}(\zeta_5)$ induziert auf den rationalen Basen lediglich triviale Galois-Effekte, da die Zyklizität des Frobenius-Automorphismus modulo $5$ die geometrische Starrheit des Raums unberührt lässt.
Wir betrachten daher den nichttrivialen algebraischen Lift auf die Hurwitz-Ordnung $\mathcal{H}$ des quaternionischen Gitters. Die elementaren Zustände $E, A, B, C$ werden hierbei mit der kanonischen Basis $\{1, i, j, k\}$ identifiziert. Ein allgemeiner lokaler Zustand $q \in \mathcal{S}_{\mathrm{klein}}$ formalisiert sich als:
\begin{equation}
    q = e \cdot 1 + a \cdot i + b \cdot j + c \cdot k \in \mathcal{H}
\end{equation}
Der über den geschlossenen Rundweg $\Gamma = R \cdot K$ definierte Mittelungsoperator
\begin{equation}
    \Pi_\Gamma = \frac{1}{4} \sum_{k=0}^{3} \Gamma^k = \begin{pmatrix} 1 & 0 & 0 & 0 \\ 0 & \frac{1}{2} & 0 & \frac{1}{2} \\ 0 & 0 & 1 & 0 \\ 0 & \frac{1}{2} & 0 & \frac{1}{2} \end{pmatrix}
\end{equation}
zerlegt den Raum $\mathbb{Q}^4$ in zwei disjunkte Eigenräume. Numerische Läufe (\texttt{run\_rundweg.sh}) bestätigen die rationale Eigenraumstruktur:
\begin{align}
    \mathcal{V}_{\lambda=1} &= \operatorname{span}_{\mathbb{Q}} \left\{ (1,0,0,0)^T, (0,0,1,0)^T, (0,1,0,1)^T \right\} \\
    \mathcal{V}_{\lambda=0} &= \operatorname{span}_{\mathbb{Q}} \left\{ (0,1,0,-1)^T \right\}
\end{align}
Der Kern des Operators isoliert somit exakt den asymmetrischen Fluktuationsträger $q_{\mathrm{asym}} = i - k$, welcher der reinen Differenzkomponente $A - C$ entspricht.
%%%%%%%%%%%%%%%%%%%%%%%%%%%%%%%%%%%%%%%%%%%%%%%%%%%%%%%%%%%%%%%%%%
% KAPITEL 8: KOLLAPSPRINZIP UND QUADRATISCHER SCHWUNG
%%%%%%%%%%%%%%%%%%%%%%%%%%%%%%%%%%%%%%%%%%%%%%%%%%%%%%%%%%%%%%%%%%
\section{Das Kollapsprinzip und die arithmetische Resonanz}
\label{sec:kollapsprinzip}
Die numerische Untersuchung des Bamberger Modells für den zeitlichen Resonanzpunkt $M=113160$ zeigt, dass eine lineare Absorption des asymmetrischen Zustands $q_{\mathrm{asym}}$ innerhalb der reinen Linksideale $I_5 = \mathcal{H}\pi_5$ und $I_7 = \mathcal{H}\pi_7$ nicht stattfindet. Die rein lineare Teilbarkeit erweist sich als unzureichend, um das Zusammenspiel zwischen den diskreten Freiheitsgraden des Kerns und den kontinuierlichen globalen Zuständen abzubilden.
Zur Auflösung dieses strukturellen Engpasses wird der Begriff der \textit{arithmetischen Resonanz} über den bilinearen \textit{quadratischen Schwung} $u$ erweitert.
\begin{definition}[Quadratischer Schwung]
Sei $q \in \mathcal{H}$ ein Zustand des Kleinschen Kerns und $q' \in \mathcal{H}$ ein zugeordneter Referenz- oder Nachbarzustand der Schale $\mathcal{S}_{\mathrm{halb}}$. Der quadratische Schwung $u$ ist definiert als das quaternionische Produkt:
\begin{equation}
    u = \bar{q} \cdot q'
\end{equation}
\end{definition}
Ein arithmetischer Resonanzpunkt liegt vor, wenn nicht der Zustand selbst, sondern sein evolutionärer Schwung $u$ mit den maximalen Idealen der Primquaternionen interferiert. 
\begin{theorem}[Erweitertes Kollaps-Theorem]
\label{thm:kollaps-erweitert}
Der Kollaps des EABC-Zustandsraums vollzieht sich entlang der Kette $4 \rightarrow 2 \rightarrow 1$. Die Reduktion der lokalen Fluktionen $q_{\mathrm{asym}} \in \mathcal{V}_{\lambda=0}$ auf die glatte Schicht $\mathcal{S}_{\mathrm{glatt}}$ wird durch die Bedingung induziert, dass der quadratische Schwung $u$ der invarianten Paare eine Normabsenkung modulo der halbzähligen Hurwitz-Primquaternionen erfährt:
\begin{equation}
    N(u) \equiv 0 \pmod{p} \quad \text{für } p \in \langle 5, 7 \rangle
\end{equation}
\end{theorem}
Durch diesen Mechanismus werden die nichtinvarianten, zyklischen Komponenten destruktiv interferiert, während die makroskopische EABC-Bernoulli-Dynamik gegen den glatten Fixpunkt $\mathbf{p}_* = \frac{1}{4}(1,1,1,1)$ konvergiert.

\subsection{Numerische Verifikation}
\label{subsec:numerische-verifikation-schalen}

\begin{proposition}[Befunde der Sage-Läufe]
\label{prop:numverifikation-schalen}
Die Skripte \texttt{run\_hurwitz\_resonanz.sh}, \texttt{run\_schwung\_resonanz.sh} und \texttt{masse\_konstitution\_test.sage} bestätigen die $\mathbb{Q}$-Spektralzerlegung und die Hurwitz-Normkriterien, stützen aber für Satz~\ref{thm:kollaps-erweitert} keinen vollständigen numerischen Nachweis am Bamberger-Anker $M=113160$:
\begin{enumerate}
    \item \textbf{Linear.} Für $q_{\mathrm{asym}}=i-k$ liegt kein nichttriviales rationales Vielfaches in $I_5 \cup I_7$ (lineare $I_5/I_7$-Resonanz negativ).
    \item \textbf{Quadratisch, Referenzpaare.} Für $u=\bar{q}\,q'$ an Testpaaren $(q_{\mathrm{lokal}}, q_{\mathrm{glatt}})$ u.\,ä.\ gilt teils $N(u)\equiv 0 \pmod{5}$ (z.\,B.\ $N(u)\in\{885,30,60\}$), jedoch nie $\pmod{7}$ und nie $u\in I_5 \cup I_7$.
    \item \textbf{Quadratisch, $M=113160$.} Über den Zeitwürfel-Schwung-Scan: kein Treffer mit $N(u)\equiv 0 \pmod{5,7}$ und kein Ideal-Treffer in $I_5 \cup I_7$.
    \item \textbf{Masse-Konstitution (orthogonal).} Drei Kandidaten mit $q_{\mathrm{sicht}}=\Pi_\Gamma q\neq 0$ und $N(q_{\mathrm{sicht}})\ge 1$; das belegt sichtbare $\lambda=1$-Anteile, nicht die Ideal-Resonanz von $q_{\mathrm{asym}}$.
\end{enumerate}
\end{proposition}

\begin{remark}
Satz~\ref{thm:kollaps-erweitert} formuliert die Resonanz über $N(u)$ als strukturelle Erweiterung der linearen Ideal-Teilbarkeit. Die Null-Treffer bei $M=113160$ und die fehlende Ideal-Mitgliedschaft der getesteten Schwünge verbieten die Behauptung, der arithmetische Schalen-Kollaps sei in den Bamberger-Daten bereits realisiert; sie rechtfertigen höchstens die weitere Suche nach Paaren $(q,q')$ oder Einheitenklassen, für die $u\in I_p$ gilt.
\end{remark}
